% =============================================================
% Section 4: Fugue Architecture. This section follows the
% 2026-09-03 user framework: data mapping, attention over the
% mapped representation, and PIM microarchitecture. The
% execution policy is separated into Section 5.
% =============================================================

\section{Fugue Architecture}
\label{sec:design}

\Fugue{} is a heterogeneous GPU--PIM hierarchical memory
system for serving shared KV.
The GPU runs most of the linear layers, and the PIM stacks
hold the KV cache in their DRAM.
Each agent still appends to a logical KV cache in its own
context order.
\Fugue{} keeps reused chunks in a shared physical store and
stores the repair tokens for those chunks separately.

Three mechanisms build the store and serve shared KV in
memory.
\textbf{Sharing-aware data mapping.}
The mapping keeps shared KV dense and packs per-agent diffs.
\textbf{Attention computation in PIM.}
The attention path reads the mapped store directly.
\textbf{PIM microarchitecture.}
The hardware adds multi-query execution and the state needed
to join master and diff.

\subsection{Sharing-Aware Data Mapping}
\label{sec:layout-dim}

To keep shared chunks dense while giving each agent a
repaired context, the mapping represents one agent's logical
KV cache with two physical parts.
The master side holds the chunks reused across agents.
The diff side holds the tokens that the upstream sharing
policy recomputes for one agent, packed compactly for that
agent's reuse step.
The metadata maps each packed diff token back to its logical
position, so later attention can read a continuous logical
cache even when the reused chunks were produced earlier.
A placement table records where each master chunk is stored,
and per-agent metadata records which logical positions a
diff segment overrides.

\begin{figure}[t]
  \centering
  \includegraphics[width=\columnwidth]{fig/ver1/fig4channelsplit.pdf}
  \figexpl{(a)~Each agent appends its KV to the cache in
  turn. (b)~The \Fugue{} physical layout: a head's shared
  chunks spread over the channels that head owns by the
  placement table, and that head's recomputed tokens packed
  into diff rows on one of those same channels.}
  \caption{The KV cache position and the \Fugue{} physical
  layout of master rows, per-head diff rows, and the
  placement table.}
  \label{fig:chanlayout}
  \Description{Two panels: each agent's logical KV cache
  positions, and the \Fugue{} physical layout with a head's
  shared chunks spread over its channels by the placement
  table and its recomputed tokens packed into diff rows on
  one of those channels.}
  % TODO(2026-09-05): redraw fig/ver1/fig4channelsplit.pdf -- panel (b)
  % still shows a dedicated diff channel; the layout now keeps the diff
  % rows inside the head's own channels (design A4c/A4e).
\end{figure}

\subsubsection{Master and Per-Agent Diff}
\label{sec:provision-ratio}

A stack's $16$ channels are divided among the KV heads it
holds, and a head's KV cache is split across the channels
that head owns.
The shared KV is stored in $256$-token chunks; each chunk
spans all the banks of one channel and occupies one row of
that channel, so one activation opens one chunk.
No channel is set aside for diffs.
A diff is small, $k$ tokens of a $256$-token chunk, so the
recomputed tokens of one agent for one head are packed
contiguously into dedicated diff rows on one of that head's
own channels, alongside master rows, with their metadata on
the logic die.
Two things follow.
The master rows keep every channel the head owns, so a sweep
that reads no diff, the common case for the shared corpus,
pays nothing for the diff rows.
And because an agent's diffs for a head are packed together,
the repairs of different rounds share rows although the KV
that agent wrote in between separates them in the append
stream.
Diffs are not merged across heads: each head's diff row
holds only that head's recomputed keys and values, since the
MAC of one head's key against its own query cannot be shared
with another head, and merging would only serialize the
heads' column reads onto one channel.

\subsubsection{Conflict-Aware Channel Placement}
\label{sec:table}

Agents can select different chunks from outputs produced by
earlier agents.
Append-order storage rotates a head's chunks over the
channels that head owns in the order they are written, so
whether two chunks read together by a later request land in
the same channel is a matter of where the rotation happened
to be, and two chunks written a stripe apart always collide.
If those chunks sit in one channel, their row scans serialize
on that channel and the sweep loses channel-level
parallelism.
\Fugue{} therefore uses the shared chunk as the placement
unit and keeps chunks that are read in the same sweep on
different channels of the head whenever enough channels are
available.

A software placement table records the channel and row of
each chunk.
Every prefill and decode names the chunks it will read, so
the driver knows which chunks can be read together before the
driver places a new chunk.
The driver appends the new row to a channel of the head that
is not already used by the rows read alongside the new chunk,
if a free channel exists, otherwise to the channel those rows
use least, and then records the chosen row.
The choice is made once, when the chunk is written, and
holds for every later sweep.
During attention, each sweep consults the table and issues
the scans to the recorded channels
(Figure~\ref{fig:chanlayout}b).
The table is managed in software.

\subsubsection{Agent Attachment and Diff Metadata}
\label{sec:diff-metadata}

An agent attaches to chunks that are already stored, and the
upstream policy names the tokens it recomputes in those
chunks.
The agent's prefill computes that set together with the
new tokens the agent adds.
The two kinds are written to different places.
The recomputed tokens go to the head's diff rows as this
agent's segment and are packed contiguously.
The added tokens go to master rows as any other new content,
in the channel the placement table assigns to them.
A chunk reused whole adds nothing to the diff rows.

The master and the diff are joined logically on the logic
die.
The per-agent metadata gives each compacted diff token's
logical position, the context slot that token overrides
for this agent.
The driver keeps one entry per agent and token beside the
segment and loads those entries when the agent attaches.
An agent that leaves frees its segment, and the hole stays
until every sweep of that channel has finished.
Compaction closes the hole after those sweeps finish.

\subsection{Attention Computation in PIM}
\label{sec:attention-pim}

To compute attention on the mapped representation, \Fugue{}
keeps the master bytes fixed and changes only the query and
the merge path around the master bytes.


\subsubsection{Position-Decoupled Query--Key Computation}
\label{sec:c1}
\label{sec:c1-inv}
\label{sec:c1-where}

Reuse needs each shared key at an agent's own position, but
RoPE has encoded the producing position into the stored
key.
\Fugue{} decomposes that position into two rotations.
The stored key carries the token's relative rotation inside
its chunk.
The query variant carries the query position minus the
chunk's base position in the agent's context.
The GPU produces one query variant per chunk scan because it
already produces the query and holds each agent's segment
offsets.
The link carries the rotated variants, as a decode kernel on
a GPU does within one device~\cite{lazyattn}.
The same stored copy can answer every agent.
Diff tokens are written under the same alignment, so the two
sides of the store meet at the same token-wise positions.

\subsubsection{Master--Diff Score and Context Merge}
\label{sec:diemerge}

Write $\Qprime_i$ for rotated query $i$, $k_t$ for the
master key at position $t$, $\tilde{k}_{i,t}$ for the key
recomputed at $t$ in query $i$'s context, and
$\mathcal{D}_i$ for the positions recomputed in that
context.
On the way in, the die assembles one score per position of
query $i$'s context,
\begin{equation}
s_{i,t}=
\begin{cases}
\Qprime_i{}^{\!\top}\tilde{k}_{i,t}, & t\in\mathcal{D}_i,\\
\Qprime_i{}^{\!\top}k_t, & t\notin\mathcal{D}_i,
\end{cases}
\label{eq:step2}
\end{equation}
The master channels supply the scores for positions outside
$\mathcal{D}_i$, and the diff rows supply the scores for
the recomputed positions.

On the way back, with $T_j$ the token share of bank group
$j$ and $p_i=\mathrm{softmax}(s_i)$, the die sends each
probability to exactly one side.
The routing rule is
\begin{equation}
\mathrm{dst}_{i,t}=
\begin{cases}
\mathrm{diff}, & t\in\mathcal{D}_i,\\
\mathrm{master}\ j, & t\in T_j \land t\notin\mathcal{D}_i.
\end{cases}
\label{eq:step4}
\end{equation}
The die sends $p_{i,t}$ to $\mathrm{dst}_{i,t}$.
Every token's score and context contribution therefore comes
from exactly one side.
The merged output is the dense computation, term by term and
up to the order of the sums.

\subsubsection{Multi-Query Shared-Row Dataflow}
\label{sec:mq-dataflow}

The dataflow has one invariant: a row opened for shared K or
V is used by all active queries before the next row opens.
Append attention supplies several active queries from one
agent during prefill.
Decode supplies one active query from each of several agents.
Both cases create multiple active queries for the same opened
row, so \Fugue{} needs multi-query execution inside the bank.
The source of the active queries changes, but the row sweep
does not.

\begin{figure}[t]
  \centering
  \includegraphics[width=\columnwidth]{fig/ver1/fig6Memoryflow.pdf}
  \figexpl{(a)~The three levels and the six operand flows:
  the query and the probabilities descend, the scores and the
  partial contexts ascend. (b)~$K^{\!\top}$ split across
  channel, pseudo-channel, bank group, and bank. (c)~$V$
  split across the same hierarchy, transposed. (d)~One
  \texttt{MAC\_AB} command opens a row for all active
  queries before the next row opens.}
  \caption{Dataflow of one head through the memory system.}
  \label{fig:memflow}
  \Description{Four panels: a three-level hierarchy with six
  operand flows for the descending query and probabilities
  and the ascending scores and contexts, the $K^{\!\top}$ and
  $V$ splits across channel, pseudo-channel, bank group, and
  bank, and one command opening a row for all active queries
  before the next row opens.}
\end{figure}

The score and context passes use the same master--diff
routing as a single query.
The difference is that one opened row must feed several
query-side or probability-side operands.
This extra work changes the per-column arithmetic, not the
sequence of rows the sweep reads.

\subsection{PIM Microarchitecture}
\label{sec:microarch}
\label{sec:die}
\label{sec:layout}

The microarchitecture spreads the attention across three
levels: the bank, the bank group, and the logic die
(Figure~\ref{fig:microarch}).

\begin{figure*}[t]
  \centering
  \includegraphics[width=\textwidth]{fig/ver1/fig5Microarch.pdf}
  \figexpl{(a)~The two per-bank GEMVs, the score
  $Q\,K^{\!\top}$ and the context $P\,V$; the $B$ resident
  queries share every scanned column of the key slice, while
  a probability entry is used only
  $\lceil(\dhead/4)/16\rceil{=}2$ times. (b)~The bank with
  its row buffer, pipelined GEMV unit, and double-buffered
  GEMV buffer of resident queries; the bank group with its
  accumulator; and the per-channel die units, the
  accumulator, the $\mathcal{D}_i$ map, and the softmax unit
  with its buffers. (c)~The system: the xPUs, the
  memory-system controller with its TLB and HBM controllers,
  and the stacks.}
  \caption{The microarchitecture at three levels, the bank,
  the channel, and the system.}
  \label{fig:microarch}
  \Description{Three panels: the two per-bank GEMV shapes for
  score and context; the bank with its row buffer, GEMV unit,
  and double-buffered query buffer, the bank group with its
  accumulator, and the per-channel accumulator,
  $\mathcal{D}_i$ map, and softmax unit; and the system with
  its xPUs, memory-system controller, TLB, HBM controllers,
  and stacks.}
\end{figure*}

\subsubsection{Bank-Level MQ-MAC}
\label{sec:mq}

\Fugue{} keeps the bank-level PIM command substrate: one PE
instruction per column command in all-bank mode, with average
stack power bounded by the original HBM envelope~\cite{hbmpim,attacc}.
The new bank-level behavior is multi-query use of one column
read.
The multi-query MAC (MQ-MAC) command still moves one column
of the shared K or V from the row buffer into the PE.
\Fugue{} changes how many query-side or probability-side
operands consume that column.

The two GEMV shapes determine operand placement.
On the score side, one query slice is reused across the
columns of a key row, so queries stay resident.
On the context side, each probability entry is consumed only
by the value columns for the output slice the bank holds, so
probabilities stream over the through-silicon-via (TSV) bus.
The bank's capacity axis therefore binds only the resident
queries.

Holding the queries resident reuses the existing in-bank
GEMV buffer~\cite{attacc}.
The stock buffer of $S$ bytes, sized for the context phase,
is far from full on the score side.
One $64$-B query slice leaves room for
$n_{\mathrm{cap}}=S/64$ query slices, eight in AttAcc's
$512$-B buffer, so \Fugue{} holds that many query slices.
The queries arrive and are stored in the GEMV buffer.
As MQ-MAC reads one column out of the row buffer,
the queries multiply that column in round-robin order, and
the empty slots of the queue are jumped over.
The multiplies stream one query per cycle, leaving time for
the products to go through the adder tree and reach the
accumulator in the bank group.

On the context side the probabilities stream rather than
reside.
As MQ-MAC moves a column of V out of the opened
row, each query's probability slice for that column arrives
over the TSV bus.
The probability is used twice, against the slice of the
output dimension the bank holds
(Figure~\ref{fig:memflow}d).
The context pass is therefore bounded by the TSV bus and its
read-write turnaround rather than by the bank's capacity.

The lever is the PE clock.
With $n$ queries resident, the $n$ MACs share one column
read, so the bank does more useful work per read.
The $n$ MACs must still finish inside the column-command
interval.
AttAcc's command interval is $6$ cycles ($t_{CK}$) for its
single MAC.
Under \Fugue{}, MQ-MAC carries $n$ operands for the same
column.
With $f_{\mathrm{PE}}$ the PE clock,
$E_{\mathrm{col}}$ a column read, $E_{\mathrm{op}}$ one MAC,
and $E_6$ the energy budget of AttAcc's $6$-cycle command,
the interval in whole cycles is
\begin{equation}
I(n,f_{\mathrm{PE}})=\max\!\Bigl(6,\ \bigl\lceil n/(f_{\mathrm{PE}}\,t_{CK})\bigr\rceil,\
\bigl\lceil 6\,(E_{\mathrm{col}}+n\,E_{\mathrm{op}})/E_6\bigr\rceil\Bigr),
\label{eq:interval}
\end{equation}
a compute term that falls with the clock and an energy term
constrained by the same budget.
The two meet at the balanced frequency point
$f_{\mathrm{PE}}^\star$.
Past $f_{\mathrm{PE}}^\star$, the interval stays flat.
At $f_{\mathrm{PE}}^\star$, the single-stage multiply-add
path meets timing, so the PE stays a single stage.

\subsubsection{Bank-Group Accumulation}
\label{sec:bankgroup}

The bank group joins what its four banks produce.
The four banks hold complementary slices of the head.
On the key side, each bank carries a $\dhead/4$ slice of
every key, and an adder in the group reduces the four
same-column partial products before they enter the
per-query accumulator.
On the value side, each bank carries a slice of the output
dimension, and the group passes its partials
through~\cite{attacc}.
With $n_{\mathrm{cap}}$ queries resident, the group keeps one
partial score and one partial context per query in its
accumulator.
The per-query accumulator is the only growth at this level.

\subsubsection{Logic-Die Merge and Softmax}
\label{sec:logicdie}

The logic die holds the state that makes the compact
master--diff store look dense to each query.
It runs the softmax and realizes the merge from the metadata
of each context in use, which the driver loads.
A decoder holds each compacted diff token's logical
position, and those positions form $\mathcal{D}_i$.
The $\mathcal{D}_i$ bitmap filters the master banks' writes
into the softmax buffer, so the result is the same whichever
side arrives first.
On the way back, the same bitmap drives the mask gate on the
recirculation path.
An adder then sums the partial contexts across the banks,
the channels, and the two sides into the per-head context
that leaves the stack~\cite{attacc}.

\subsubsection{Hardware and Command Compatibility}
\label{sec:compat}

The bank PE keeps its MAC as the only arithmetic, and the
DRAM command set stays unchanged.
The additions span the three levels.
The bank adds the resident-query buffer and per-query state.
The group adds the per-query accumulator.
The die adds the decoder, the $\mathcal{D}_i$ bitmap with
its write filter and mask gate, a causal comparator, and the
context adder.
